% answers.tex —— 1990 年上海卷数学 · 参考答案与详细解析片段
% 由 tools/generate_projects.py 依来源总账自 1990/1990_数学真题_上海_试题.md + 1990/1990_数学真题_上海_解析正文.md + 1990/1990_数学真题_上海_答案要点.md 生成，勿手改（改源 MD 后重跑）


\begin{gktip}
以下推导依据本项目转录稿独立完成，已完成整卷独立盲解与差异裁定，但不是来源印载解析。第 18 题保留原文记号冲突；第 23 题的推导以“点 $P$ 在直线 $x=2$ 上移动”为条件性释读，不据此改写原题。
\end{gktip}

\begin{answer}
\gksolhead{1}{}

\textbf{题目}：

函数 $y=\dfrac{\sqrt{x+4}}{x+2}$ 的定义域是\gkblank{}。

\Needspace{4\baselineskip}
\textbf{解析}

平方根要求 $x+4\ge0$，故 $x\ge-4$；分母还要求 $x+2\ne0$。合并两项限制即可。

\gkanswer{答案：$[-4,+\infty)\setminus\{-2\}$}

\end{answer}
\begin{answer}
\gksolhead{2}{}

\textbf{题目}：

函数 $y=\arcsin x$（$x\in[-1,1]$）的反函数是\gkblank{}。

\Needspace{4\baselineskip}
\textbf{解析}

$y=\arcsin x$ 将 $[-1,1]$ 一一映到 $[-\frac\pi2,\frac\pi2]$。交换变量后，反函数为 $y=\sin x$，定义域为 $[-\frac\pi2,\frac\pi2]$。

\gkanswer{答案：$y=\sin x$，$x\in[-\frac\pi2,\frac\pi2]$}

\end{answer}
\begin{answer}
\gksolhead{3}{}

\textbf{题目}：

过点 $(1,2)$ 且与直线 $2x+y-1=0$ 平行的直线方程是\gkblank{}。

\Needspace{4\baselineskip}
\textbf{解析}

所求直线与 $2x+y-1=0$ 平行，法向量相同，故可写为 $2x+y+c=0$。代入 $(1,2)$ 得 $c=-4$。

\gkanswer{答案：$2x+y-4=0$}

\end{answer}
\begin{answer}
\gksolhead{4}{}

\textbf{题目}：

已知圆柱的轴截面是正方形，它的面积是 $4\,\mathrm{cm}^2$，那么这个圆柱的体积是\gkblank{}$\,\mathrm{cm}^3$。（结果中保留 $\pi$）

\Needspace{4\baselineskip}
\textbf{解析}

设圆柱底面半径为 $r$、高为 $h$。轴截面为正方形意味着 $h=2r$，且其面积 $(2r)h=4$，所以 $r=1$、$h=2$。圆柱体积为 $\pi r^2h$。

\gkanswer{答案：$2\pi\,\mathrm{cm}^3$}

\end{answer}
\begin{answer}
\gksolhead{5}{}

\textbf{题目}：

在 $\triangle ABC$ 中，已知 $\cos A=-\dfrac35$，则 $\sin\dfrac A2=$\gkblank{}。

\Needspace{4\baselineskip}
\textbf{解析}

三角形内角 $A\in(0,\pi)$，故半角正弦为正。由半角公式，$\sin\frac A2=\sqrt{\frac{1-\cos A}{2}}=\sqrt{\frac{1+3/5}{2}}$。

\gkanswer{答案：$\frac{2}{\sqrt5}$}

\end{answer}
\begin{answer}
\gksolhead{6}{}

\textbf{题目}：

设复数 $\omega=\cos\dfrac{2\pi}{5}+i\sin\dfrac{2\pi}{5}$，则 $\omega+\omega^2+\omega^3+\omega^4+\omega^5$ 的值是\gkblank{}。

\Needspace{4\baselineskip}
\textbf{解析}

$\omega$ 是满足 $\omega^5=1$ 的五次单位根，且 $\omega\ne1$。由等比数列求和，$1+\omega+\omega^2+\omega^3+\omega^4=0$；所以题设五项和为 $(1+\omega+\cdots+\omega^4)-1+\omega^5=0$。

\gkanswer{答案：$0$}

\end{answer}
\begin{answer}
\gksolhead{7}{}

\textbf{题目}：

已知圆锥的中截面周长为 $a$，母线长为 $l$，则它的侧面积等于\gkblank{}。

\Needspace{4\baselineskip}
\textbf{解析}

设底面半径为 $R$。圆锥的平行底面截面与底面相似；中截面半径为 $R/2$，所以其中截面周长 $a=\pi R$。侧面积 $S=\pi Rl$，代入 $\pi R=a$ 得所求关系。

\gkanswer{答案：$al$}

\end{answer}
\begin{answer}
\gksolhead{8}{}

\textbf{题目}：

已知 $(x+a)^7$ 的展开式中，$x^4$ 的系数是 $-280$，则实数 $a=$\gkblank{}。

\Needspace{4\baselineskip}
\textbf{解析}

$(x+a)^7$ 中 $x^4$ 项系数为 $\binom74a^3=35a^3$。令其等于 $-280$，得到 $a^3=-8$。

\gkanswer{答案：$a=-2$}

\end{answer}
\begin{answer}
\gksolhead{9}{}

\textbf{题目}：

双曲线 $2mx^2-my^2=2$ 的一条准线是 $y=1$，则 $m=$\gkblank{}。

\Needspace{4\baselineskip}
\textbf{解析}

$m>0$ 时双曲线的横向项为正，准线平行于 $x$ 轴，不可能是 $y=1$。故 $m<0$，改写为 $\frac{y^2}{-2/m}-\frac{x^2}{-1/m}=1$。此时 $a^2=-2/m$、$b^2=-1/m$、$c^2=a^2+b^2=-3/m$，准线为 $y=\pm a/e$，其中 $e=c/a=\sqrt{3/2}$。由一条准线为 $y=1$ 得 $a/e=1$，解得 $m=-\frac43$。

\gkanswer{答案：$m=-\frac43$}

\end{answer}
\begin{answer}
\gksolhead{10}{}

\textbf{题目}：

平面上，四条平行直线与另外五条平行直线互相垂直，则它的矩形共有\gkblank{}个。（结果用数值表示）

\Needspace{4\baselineskip}
\textbf{解析}

每个矩形由第一组平行线中任选两条、第二组平行线中任选两条唯一确定，故数量为 $\binom42\binom52$。

\gkanswer{答案：$60$}

\end{answer}
\begin{answer}
\gksolhead{11}{}

\textbf{题目}：

圆的半径是 $1$，圆心的极坐标是 $(1,0)$，则这个圆的极坐标方程是\mbox{（　　）}。

（A）$\rho=\cos\theta$　　（B）$\rho=\sin\theta$　　（C）$\rho=2\cos\theta$　　（D）$\rho=2\sin\theta$

\Needspace{4\baselineskip}
\textbf{解析}

圆心为笛卡尔坐标 $(1,0)$，半径为 $1$。代入极坐标 $x=\rho\cos\theta$、$y=\rho\sin\theta$：$(\rho\cos\theta-1)^2+\rho^2\sin^2\theta=1$，化简为 $\rho^2-2\rho\cos\theta=0$。圆的极坐标方程可写为 $\rho=2\cos\theta$。

\gkanswer{答案：C}

\end{answer}
\begin{answer}
\gksolhead{12}{}

\textbf{题目}：

函数 $f(x)$ 和 $g(x)$ 的定义域均为 $\mathbb R$，“$f(x),g(x)$ 都是奇函数”是“$f(x)$ 与 $g(x)$ 的积是偶函数”的\mbox{（　　）}。

（A）必要条件但非充分条件　　（B）充分条件但非必要条件

（C）充分必要条件　　（D）非充分条件也非必要条件

\Needspace{4\baselineskip}
\textbf{解析}

若 $f,g$ 都是奇函数，则 $f(-x)g(-x)=(-f(x))(-g(x))=f(x)g(x)$，所以乘积是偶函数。反向不成立，例如取 $f(x)=g(x)=1$，其乘积为偶函数，而两个因子均非奇函数。

\gkanswer{答案：B}

\end{answer}
\begin{answer}
\gksolhead{13}{}

\textbf{题目}：

设点 $P$ 在有向线段 $\overrightarrow{AB}$ 的延长线上，$P$ 分 $\overrightarrow{AB}$ 所成的比为 $\lambda$，则\mbox{（　　）}。

（A）$\lambda<-1$　　（B）$-1<\lambda<0$　　（C）$0<\lambda<1$　　（D）$\lambda>1$

\Needspace{4\baselineskip}
\textbf{解析}

取有向直线坐标 $A=0,B=1$。延长线段 $AB$ 越过 $B$ 上的点可写成 $P=t$ 且 $t>1$。有向分比 $\lambda=\frac{\overline{AP}}{\overline{PB}}=\frac{t}{1-t}<-1$。

\gkanswer{答案：A}

\end{answer}
\begin{answer}
\gksolhead{14}{}

\textbf{题目}：

设 $2^a=3$，$2^b=6$，$2^c=12$，则数列 $a,b,c$\mbox{（　　）}。

（A）是等差数列但不是等比数列　　（B）是等比数列但不是等差数列

（C）既是等差数列又是等比数列　　（D）既不是等差数列又不是等比数列

\Needspace{4\baselineskip}
\textbf{解析}

$a=\log_2 3$，$b=\log_2 6=a+1$，$c=\log_2 12=a+2$，三项成等差数列。等比数列还需 $b^2=ac$；但 $ac=(a)(a+2)=a^2+2a$，$b^2=(a+1)^2=a^2+2a+1$，二者不等。

\gkanswer{答案：A}

\end{answer}
\begin{answer}
\gksolhead{15}{}

\textbf{题目}：

设 $\alpha$ 角属于第二象限，且 $\left|\cos\dfrac\alpha2\right|=-\cos\dfrac\alpha2$，则 $\dfrac\alpha2$ 角属于\mbox{（　　）}。

（A）第一象限　　（B）第二象限　　（C）第三象限　　（D）第四象限

\Needspace{4\baselineskip}
\textbf{解析}

按任意角的标准定义，“$\alpha$ 属于第二象限”表示其终边在第二象限，故可写为 $\alpha=\alpha_0+2k\pi$，其中 $\alpha_0\in(\frac\pi2,\pi)$、$k\in\mathbb Z$。于是 $\alpha/2=\alpha_0/2+k\pi$，且 $\alpha_0/2\in(\frac\pi4,\frac\pi2)$。条件 $|\cos(\alpha/2)|=-\cos(\alpha/2)$ 等价于 $\cos(\alpha/2)\le0$，这恰在 $k$ 为奇数时成立；此时 $\alpha/2$ 的终边在第三象限，故选（C）。这一解读也与一份明确列出本题的三角函数专题资料所给答案及其 $k\in\mathbb Z$ 解释一致；该资料是外部交叉佐证，不是本卷官方答案来源。

\gkanswer{答案：C}

\end{answer}
\begin{answer}
\gksolhead{16}{}

\textbf{题目}：

设过长方体同一个顶点的三个面的对角线长分别是 $a,b,c$，那么这个长方体的对角线长是\mbox{（　　）}。

（A）$\sqrt{a^2+b^2+c^2}$　　（B）$\sqrt{\dfrac{a^2+b^2+c^2}{2}}$

（C）$\sqrt{\dfrac{a^2+b^2+c^2}{3}}$　　（D）$\dfrac{\sqrt{a^2+b^2+c^2}}{2}$

\Needspace{4\baselineskip}
\textbf{解析}

设过同一顶点的三条棱长为 $u,v,w$。三个面上的对角线平方分别是 $u^2+v^2$、$u^2+w^2$、$v^2+w^2$，其和为 $2(u^2+v^2+w^2)$。长方体体对角线平方为 $u^2+v^2+w^2$，故等于 $\frac{a^2+b^2+c^2}{2}$。

\gkanswer{答案：B}

\end{answer}
\begin{answer}
\gksolhead{17}{}

\textbf{题目}：

函数 $f(x)=a\tan\dfrac xa$ 的最小正周期是\mbox{（　　）}。

（A）$\pi a$　　（B）$\pi|a|$　　（C）$\dfrac\pi a$　　（D）$\dfrac\pi{|a|}$

\Needspace{4\baselineskip}
\textbf{解析}

$a\ne0$ 时，$f(x+T)=f(x)$ 等价于 $\tan(\frac{x+T}{a})=\tan(\frac xa)$，最小正的自变量平移满足 $T/|a|=\pi$。故最小正周期为 $\pi|a|$。

\gkanswer{答案：B}

\end{answer}
\begin{answer}
\gksolhead{18}{}

\textbf{题目}：

已知 $1<x<d$，令 $a=(\log_d x)^2$，$b=\log_d(x^2)$，$b=\log_d(\log_d x)$，则\mbox{（　　）}。

（A）$a<b<c$　　（B）$a<c<b$　　（C）$c<b<a$　　（D）$c<a<b$

\Needspace{4\baselineskip}
\textbf{解析}

原文先定义 $b=\log_d(x^2)$，随后又定义 $b=\log_d(\log_d x)$，符号冲突使三个量及待比较关系无法按字面同时成立。若只把第三个量记号假设为 $c$，则由 $1<x<d$ 可知 $d>1$ 且 $t=\log_d x\in(0,1)$。于是 $a=t^2$、$b=2t$、$c=\log_d t<0$，故 $c<a<b$。此为条件性推论，不能替代原题校核。

\gkanswer{答案：U/PENDING；原文两次定义 $b$；若将第三个量视为 $c$，则条件性答案为 D：$c<a<b$。}

\end{answer}
\begin{answer}
\gksolhead{19}{}

\textbf{题目}：

设 $a,b$ 是两条异面直线，那么下列四个命题中的假命题是\mbox{（　　）}。

（A）经过直线 $a$ 有且只有一个平面平行于直线 $b$

（B）经过直线 $a$ 有且只有一个平面垂直于直线 $b$

（C）存在分别经过直线 $a$ 和 $b$ 的两个互相平行的平面

（D）存在分别经过直线 $a$ 和 $b$ 的两个互相垂直的平面

\Needspace{4\baselineskip}
\textbf{解析}

设两异面直线方向分别为 $\mathbf u,\mathbf v$。存在包含 $a$ 且平行于 $b$ 的平面；但若 $a$ 与 $b$ 不垂直，任何包含 $a$ 的平面都不能以 $b$ 为法向量，因此不存在垂直于 $b$ 且包含 $a$ 的平面。题项（B）并非对任意两条异面直线都成立，是假命题。两异面直线的方向不平行，取 $\mathbf n_1=\mathbf u\times\mathbf v$ 和 $\mathbf n_2=\mathbf u-\frac{\mathbf u\cdot\mathbf v}{\mathbf v\cdot\mathbf v}\mathbf v$，二者分别垂直于 $a,b$ 的方向且互相垂直，因此分别以这两向量作法向可构造互相垂直、分别包含 $a,b$ 的平面；这也说明（D）可成立。将一个含 $a$ 的平面平移至过 $b$ 的平行平面则说明（C）可成立。

\gkanswer{答案：B}

\end{answer}
\begin{answer}
\gksolhead{20}{}

\textbf{题目}：

下列四个函数中，在定义域内不具有单调性的函数是\mbox{（　　）}。

（A）$y=\cot(\arccos x)$　　（B）$y=\tan(\arcsin x)$

（C）$y=\sin(\arctan x)$　　（D）$y=\cos(\arctan x)$

\Needspace{4\baselineskip}
\textbf{解析}

前三个函数分别化为 $\frac{x}{\sqrt{1-x^2}}$、$\frac{x}{\sqrt{1-x^2}}$、$\frac{x}{\sqrt{1+x^2}}$（各自定义域内），均严格递增。第四个函数为 $\frac1{\sqrt{1+x^2}}$，是偶函数，在整个 $\mathbb R$ 上不单调。

\gkanswer{答案：D}

\end{answer}
\begin{answer}
\gksolhead{21}{}

\textbf{题目}：

已知 $\log_5(x^2+2x-2)=0$，$2\log_5(x+2)-\log_5 y+\dfrac12=0$，求 $y$ 的值。

\Needspace{4\baselineskip}
\textbf{解析}

第一式给 $x^2+2x-2=1$，所以 $x=1$ 或 $x=-3$；第二式中的 $\log_5(x+2)$ 要求 $x>-2$，故只取 $x=1$。代入第二式得 $\log_5 9-\log_5y=-\frac12$，即 $\log_5(9/y)=-\frac12$，因此 $y=9\sqrt5$。

\gkanswer{答案：$y=9\sqrt5$}

\end{answer}
\begin{answer}
\gksolhead{22}{}

\textbf{题目}：

求方程 $\sqrt{5\cos x+\cos 2x}+\sin x=0$ 在 $[0,2\pi)$ 上的解。

\Needspace{4\baselineskip}
\textbf{解析}

根式非负，原等式迫使 $\sin x\le0$。平方并用 $\cos2x=1-2\sin^2x$：$5\cos x+1-2\sin^2x=\sin^2x$，即 $3\cos^2x+5\cos x-2=0$。其根为 $\cos x=\frac13$ 或 $-2$，后者不可能。结合 $x\in[0,2\pi)$ 且 $\sin x\le0$，保留下半圆上的解 $x=2\pi-\arccos\frac13$；代回时根式等于 $|\sin x|=-\sin x$，满足原式。

\gkanswer{答案：$x=2\pi-\arccos\frac13$}

\end{answer}
\begin{answer}
\gksolhead{23}{}

\textbf{题目}：

已知点 $P$ 直线 $x=2$ 上移动，直线 $l$ 通过原点且与 $OP$ 垂直；通过点 $A(1,0)$ 及点 $P$ 的直线 $m$ 和直线 $l$ 交于点 $Q$。求点 $Q$ 的轨迹方程，并指出该轨迹的名称和它的焦点坐标。

\Needspace{4\baselineskip}
\textbf{解析}

\textbf{条件性推导：}假设原文意为点 $P$ 在直线 $x=2$ 上移动，令 $P=(2,t)$。由于 $l\perp OP$ 且过原点，$l$ 的方程为 $2x+ty=0$；直线 $AP$ 的方程为 $y=t(x-1)$。解两式得 $Q=(\frac{t^2}{t^2+2},-\frac{2t}{t^2+2})$。消去实参数 $t$，得到 $4(x-\frac12)^2+2y^2=1$。当有限 $t$ 取值时不取到 $(1,0)$，故几何轨迹是该椭圆去掉点 $(1,0)$。半轴平方为 $1/4$ 与 $1/2$，焦距满足 $c^2=1/2-1/4=1/4$，焦点为 $(\frac12,\pm\frac12)$。若原题缺字并非“在”，则本推导不能视为无条件答案。

\gkanswer{答案：在 $P$ 位于直线 $x=2$ 的条件下，轨迹为 $4(x-\frac12)^2+2y^2=1$（不含 $(1,0)$），焦点为 $(\frac12,\pm\frac12)$}

\end{answer}
\begin{answer}
\gksolhead{24}{}

\textbf{题目}：

\Needspace{5\baselineskip}
已知直线 $l:x-ny=0$（$n\in\mathbb N$）；圆 $M:(x+1)^2+(y+1)^2=1$；抛物线 $\Phi:y=(x-1)^2$。又 $L$ 与 $M$ 交于点 $A,B$；$L$ 与 $\Phi$ 交于点 $C,D$。求
\nopagebreak[4]
\[
\lim_{n\to\infty}\frac{|AB|^2}{|CD|^2}.
\]

\Needspace{4\baselineskip}
\textbf{解析}

直线可写为 $y=x/n$。圆心 $(-1,-1)$ 到该直线的距离平方为 $\frac{(n-1)^2}{n^2+1}$，故圆弦长平方 $|AB|^2=4(1-\frac{(n-1)^2}{n^2+1})=\frac{8n}{n^2+1}$。抛物线交点的横坐标满足 $x=n(x-1)^2$，其两根差的绝对值为 $\frac{\sqrt{4n+1}}n$。两交点在斜率 $1/n$ 的直线上，因此 $|CD|^2=(1+\frac1{n^2})\frac{4n+1}{n^2}=\frac{(n^2+1)(4n+1)}{n^4}$。两平方长度相除后取极限，得到 $2$。

\gkanswer{答案：$2$}

\end{answer}
\begin{answer}
\gksolhead{25}{}

\textbf{题目}：

\Needspace{5\baselineskip}
关于实数 $x$ 的不等式
\nopagebreak[4]
\[
\left|x-\frac{(a+1)^2}{2}\right|\leq\frac{(a-1)^2}{2}
\]
\Needspace{5\baselineskip}
与
\nopagebreak[4]
\[
x^2-3(a+1)x+2(3a+1)\leq0
\]
的解集依次记为 $A$ 与 $B$。求使 $A\subseteq B$ 的 $a$ 的取值范围（其中 $a\in\mathbb R$）。

\Needspace{4\baselineskip}
\textbf{解析}

绝对值不等式的解集为 $A=[2a,a^2+1]$，因为两个端点之差 $(a-1)^2\ge0$。第二个二次式判别式为 $(3a-1)^2$，两根为 $2$ 与 $3a+1$，故 $B$ 是这两根间的闭区间。

若 $a\ge\frac13$，则 $B=[2,3a+1]$。包含关系要求 $2a\ge2$ 且 $a^2+1\le3a+1$，合并得 $1\le a\le3$。若 $a<\frac13$，则 $B=[3a+1,2]$；包含关系要求 $2a\ge3a+1$ 且 $a^2+1\le2$，分别给出 $a\le-1$ 与 $-1\le a\le1$，交集仅为 $a=-1$。

\gkanswer{答案：$a\in\{-1\}\cup[1,3]$}

\end{answer}
\begin{answer}
\gksolhead{26}{}

\textbf{题目}：

如图，平面 $\alpha,\beta$ 相交于直线 $MN$，点 $A$ 在平面 $\alpha$ 上，点 $B$ 在平面 $\beta$ 上，点 $C$ 在直线 $MN$ 上，$\angle ACM=\angle BCN=45^\circ$，$A$—$MN$—$B$ 是 $60^\circ$ 的二面角，$AC=1$。求：

（1） 点 $A$ 到平面 $\beta$ 的距离；

（2） 二面角 $A$—$BC$—$M$ 的大小（用反三角函数表示）。
\nopagebreak[4]
\begin{center}% 1990 上海卷第 26 题：两平面及点线拓扑按原卷重绘，不表达精确比例。
\begin{tikzpicture}[x=0.78cm,y=0.74cm,line cap=round,line join=round]
  \coordinate (M) at (0,0);
  \coordinate (N) at (6,0);
  \coordinate (C) at (2.55,0);
  \coordinate (A) at (3.15,1.05);
  \coordinate (B) at (4.45,-0.9);
  \coordinate (AlphaLeft) at (1.25,1.55);
  \coordinate (AlphaRight) at (7.25,1.55);
  \coordinate (BetaLeft) at (1.2,-1.45);
  \coordinate (BetaRight) at (7.2,-1.45);

  % 两平面共用 MN；远侧的 alpha 边按原图用虚线表示。
  \draw[thin] (M)--(N)--(AlphaRight);
  \draw[thin,dashed] (AlphaRight)--(AlphaLeft);
  \draw[thin] (AlphaLeft)--(M);
  \draw[thin] (N)--(BetaRight)--(BetaLeft)--(M);

  % 原图中 A、B、C 的连线拓扑；透视坐标仅为排布，不代表测量值。
  \draw[thick] (A)--(C) (C)--(B) (A)--(B);

  \node[below left=2pt] at (M) {$M$};
  \node[below right=2pt] at (N) {$N$};
  \node[below=4pt] at (C) {$C$};
  \node[above left=1pt] at (A) {$A$};
  \node[below right=1pt] at (B) {$B$};
  \node[above right=3pt] at (AlphaRight) {$\alpha$};
  \node[below=5pt] at (BetaLeft) {$\beta$};
  \node[font=\scriptsize,align=center] at (3.6,-2.15)
    {按原卷点线关系重绘；图形不按比例。};
\end{tikzpicture}
\end{center}
\nopagebreak[4]
\Needspace{4\baselineskip}
\textbf{解析}

\gksubpart{（1）}

取 $C$ 为原点、$MN$ 为 $x$ 轴，令平面 $\alpha$ 为 $xy$ 平面，平面 $\beta$ 绕 $x$ 轴相对 $\alpha$ 转过 $60^\circ$。按题给角度可取 $A=(-\frac1{\sqrt2},\frac1{\sqrt2},0)$；平面 $\beta$ 的方程为 $-\sqrt3y+z=0$。点到平面的距离为 $\frac{| -\sqrt3/\sqrt2|}{\sqrt{3+1}}=\frac{\sqrt6}{4}$。

\gksubpart{（2）}

为求第二个二面角，令 $|CB|=1$（所求仅依赖其方向），并按 $\angle BCN=45^\circ$ 取 $B=(\frac1{\sqrt2},\frac1{2\sqrt2},\frac{\sqrt3}{2\sqrt2})$，再取 $M=(-1,0,0)$。沿棱 $BC$ 将 $A$、$M$ 的向量投影到 $BC$ 的垂直平面，分别记为 $A_\perp,M_\perp$。直接计算得 $|A_\perp|=\frac{\sqrt{15}}4$、$|M_\perp|=\frac1{\sqrt2}$、$A_\perp\cdot M_\perp=\frac3{4\sqrt2}$，故两半平面所成角的余弦为 $\frac{A_\perp\cdot M_\perp}{|A_\perp||M_\perp|}=\sqrt{\frac35}$，角为 $\arccos\sqrt{\frac35}=\arctan\sqrt{\frac23}$。

\gkanswer{答案（1）：$\frac{\sqrt6}{4}$}

\gkanswer{答案（2）：$\arccos\sqrt{\frac35}$}

\end{answer}
\begin{answer}
\gksolhead{27}{}

\textbf{题目}：

复平面上点 $A,B$ 对应的复数分别为 $z_1=2,z_2=-3$，点 $P$ 对应的复数为 $z$，$\dfrac{z-z_1}{z-z_2}$ 的辐角主值为 $\varphi$。当点 $P$ 在以原点为圆心、$1$ 为半径的上半圆周（不包括两个端点）上运动时，求 $\varphi$ 的最小值。

\Needspace{4\baselineskip}
\textbf{解析}

\Needspace{5\baselineskip}
令 $z=e^{it}$，$0<t<\pi$。有
\nopagebreak[4]
\[
\frac{z-2}{z+3}=\frac{(z-2)(\bar z+3)}{|z+3|^2}
=\frac{(\cos t-5)+5i\sin t}{|z+3|^2}.
\]
实部为负、虚部为正，故 $\varphi=\pi-\arctan\frac{5\sin t}{5-\cos t}$。令 $u=\cos t\in(-1,1)$，对 $\frac{5\sqrt{1-u^2}}{5-u}$ 求极值；其内部最大值在 $u=\frac15$ 取得，函数值为 $\frac5{2\sqrt6}$。该点属于开上半圆，最大值确实可取，因而 $\varphi$ 的最小值为 $\pi-\arctan\frac5{2\sqrt6}$。

\gkanswer{答案：$\pi-\arctan\frac{5}{2\sqrt6}$}

\end{answer}